\documentclass[11pt,a4paper]{article}
\usepackage{fontspec}
\setmainfont{Noto Serif}
\setsansfont{Noto Sans}
\usepackage{babel}
\babelprovide[main,import]{russian}
\usepackage{amsmath,amssymb,mathtools,booktabs,geometry,fancyhdr,microtype,xcolor,hyperref}
\geometry{left=21mm,right=21mm,top=21mm,bottom=21mm}
\definecolor{ink}{HTML}{173D57}
\hypersetup{colorlinks=true,urlcolor=ink,linkcolor=ink,pdftitle={BGS-F. Нулевая мода на степенях семи}}
\pagestyle{fancy}\fancyhf{}\lhead{\small BGS-F}\rhead{\small Нулевая мода на степенях семи}\cfoot{\thepage}
\setlength{\headheight}{15pt}\setlength{\parindent}{1em}\setlength{\parskip}{2.5pt}
\setlength{\emergencystretch}{3em}\tolerance=3000
\newcommand{\Tr}{\operatorname{Tr}}
\newcommand{\EE}{\mathbb E}
\newcommand{\R}{\mathbb R}
\newcommand{\Z}{\mathbb Z}
\newcommand{\C}{\mathbb C}
\newcommand{\si}{\mathrm{Sine}_2}
\newcommand{\eps}{\varepsilon}
\newcommand{\sinc}{\operatorname{sinc}}
\newcommand{\sect}[1]{\section*{\sffamily\color{ink}#1}}
\newcommand{\subt}[1]{\par\medskip\noindent\textbf{#1.}\quad}
\newcommand{\T}{\mathbb T}
\newcommand{\dist}{\operatorname{dist}}
\newcommand{\Ad}{\operatorname{Ad}}
\begin{document}
\noindent{\small\sffamily\color{ink} МАТЕМАТИЧЕСКАЯ ФИЗИКА\hfill 20 НОЯБРЯ 2026}
\begin{center}{\LARGE\sffamily Обычный след на степенях семи}\\[4pt]
{\large Гейзенберговские возвраты и локальный арифметический предел}
\end{center}
Пусть $U_N(t)=e^{itA_N}M_N$, где $A_N,M_N$ имеют определения исходной постановки. Физической силе $u_N$ соответствует $t=Nu_N$. В основной задаче $u_N=\tau=10^{-80}$, так что $U_N=U_N(\tau N)$.

\textbf{Фиксированный исходный кик.} Для каждой размерности $N=7^k$, $2\le k\le56$, при $m=3N/7$ получено
\[
      |\Tr U_N^m|>0.979N,\qquad
      \Delta_{2,N}(b_*)>0.040 .                                    \tag{1}
\]
Старшая размерность $7^{56}>10^{47}$. Для одного и того же ограниченного непрерывного локального функционала доказан разрыв ожиданий с $\si$ больше $0.78$ во всём диапазоне. При всех $N=7^k$, $k\ge3$, интенсивность $|\Tr U_N^m|^2$ трансцендентна для всех $m\ge1$ с $21\nmid m$, а также при скалярных линейных возвратах $m\in (3N/7)\Z_{>0}$.

\textbf{Убывающий вспомогательный кик.} Если $N=7^k\to\infty$ и $|u_N|\le\log N/(5N^2)$, локальный процесс сходится к случайно сдвинутой решётке шага $7/3$. На её точке с индексом $j$ кратность равна $2(1+\nu_7(B+j))$, где $B$ равномерен в кольце $\Z_7$, независимо от непрерывного сдвига. Для этого же режима вычислен полный предел (P), и он строго положителен. В более широком режиме $u_N=o(N^{-5/3})$ сохраняется макроскопический след при $m=3N/7$, что уже исключает $\si$ для этого режима.

На временах $m\asymp N$ с $3\nmid m$ также получена явная асимптотика обычного комплексного следа при вспомогательном параметре $|mt|\lesssim\log N/\log\log N$. После деления на линейный след остаются члены порядков ноль и один. Сумма всех высших порядков стремится к нулю за счёт арифметической компенсации. Вычислено предельное распределение коэффициента линейного отклика по времени.

\textbf{Статус.} Конечные оценки (1) относятся к неизменному $\tau$. Асимптотические теоремы о локальном процессе приведены для явно указанной убывающей силы. При $u_N=10^{-80}$ и $N\to\infty$ полный парный тест и локальный закон остаются открытыми.
\newpage
\sect{1. Линейный след и вырожденное преобразование Кэли}
Пишем $N=7^k$, $e_N(x)=\exp(2\pi i x/N)$, $\omega((a,b),(c,d))=ad-bc$ и $Q(a,b)=a^2-2b^2$. Для модульных операторов Вейля
\[
 T(d)T(e)=e_N(\bar2\omega(d,e))T(d+e),\qquad
 M T(d)M^*=T(Sd),\qquad S=\begin{pmatrix}3&4\\2&3\end{pmatrix}.
\]
Эти тождества проверены в части о возврате Пелля. Введём $S^m=\bigl(\begin{smallmatrix} a_m&2b_m\\b_m&a_m\end{smallmatrix}\bigr)$, $I_m=\operatorname{im}(S^m-I)$ на модуле $N$. Из следа сопряжения и ковариантности
\[
       |\Tr M^m|^2=|\ker(S^m-I)|\ge1.                               \tag{2}
\]
Если $v\notin I_m$, то $\Tr(T(v)M^m)=0$. Это следует сопряжением через $T(e)$ с $e\in\ker(S^m-I)$: характер $e_N(\omega(e,v))$ должен быть единицей. Симплектический аннулятор ядра равен образу по включению и сравнению мощностей. При $v=(I-S^m)e$ сопряжение $T(e)M^mT(-e)$ даёт
\[
              \frac{\Tr(T(v)M^m)}{\Tr M^m}
                         =e_N(\bar2\omega(e,S^m e)).                \tag{3}
\]
В частности, модуль отношения на носителе равен 1.

По модулю 7 матрица $S$ имеет собственные значения 2 и 4 и порядок 3. Имеем $S^3-I=14\bigl(\begin{smallmatrix}7&10\\5&7\end{smallmatrix}\bigr)$, где внутренняя матрица обратима. Биномиальная формула при возведении в степень 7 увеличивает 7-адическую кратность на один, оставляя ту же обратимую матрицу после снятия степени 7. Возведение в степень, взаимно простую с 7, умножает первый невырожденный коэффициент на единицу по модулю 7. Поэтому при $3\nmid m$ матрица $S^m-I$ обратима, а при $3\mid m$
\[
       b=\min(k,\,1+\nu_7(m)),\quad
       I_m=7^b(\Z/N\Z)^2,\quad |\Tr M^m|=7^b.                      \tag{4}
\]
Порядок $S$ равен $o=3N/7$. Для любой ненулевой исходной метки $d$ из носителя $A$ её период по модулю $7^b$, если $b>\nu_7(d)$, равен $3\,7^{b-\nu_7(d)-1}$. Здесь $\nu_7(d)$ есть минимум порядков двух координат, $\nu_7(0)=+\infty$ для нулевой координаты. Это вновь следует из обратимости оставшейся матрицы на каждой глубине. На носителе $\nu_7(d)\le1$.

Если $3\nmid m$, то из (3) и матрицы $(S^m+I)(S^m-I)^{-1}$ получаем
\[
     \frac{\Tr(T(v)M^m)}{\Tr M^m}=e_N(\kappa_m Q(v)),\qquad
     \kappa_m=-\bar4\,b_m/(a_m-1).                                  \tag{5}
\]
В знаменателе $a_m-1$ стоит единица по модулю 7, как и в $b_m$. Формула (5) верна и на общем нечётном модуле, если $\det(S^{2m}-I)$ обратим. Для проверки знака,
$(S^m+I)(S^m-I)^{-1}$ имеет нулевую диагональ и внедиагональные элементы $2b_m/(a_m-1)$, $b_m/(a_m-1)$. Подстановка в (3) даёт (5).
\newpage
\sect{2. Точный результат при фиксированном ненулевом кике}
Обозначим $p_m(t)=\Tr U_N(t)^m$, $F_m(t)=|p_m(t)|^2$ для вещественного $t$. Поскольку $A$ эрмитова и алгебраична, $p_m,F_m$ являются конечными суммами экспонент $e^{it\alpha}$ с алгебраическими вещественными частотами и алгебраическими коэффициентами. Сначала объединяем совпадающие частоты. По теореме Линдемана--Вейерштрасса в каждой ненулевой алгебраической точке $t$ значение такой суммы может быть алгебраическим только в случае её постоянства. Это следует применением теоремы к разности суммы и её значения. По (2) отсюда сразу получаем $p_m(\tau N)\ne0$ для всех нечётных $N$ и $m\ge1$. Если $F_m(\tau N)$ алгебраично, оно в точности равно $|\ker(S^m-I)|$.

Дифференцируя произведение при $t=0$ и циклически перенося множители, имеем $p'_m(0)=im\Tr(A M^m)$. Для $3\nmid m$ формула (5) даёт
\[
      \frac{F'_m(0)}m=-2\Im H(\kappa_m/N),\qquad
      H(x)=\tfrac1{4000}\sum_q C_q e^{2\pi i q x}.
                                                                     \tag{6}
\]
Здесь $C_q=4000\sum_{Q(d)=q}c_d$ для $A=\sum_d c_dT(d)$. Числа вещественны по сопряжённости $c_{-d}=\bar c_d$. Ненулевые значения $(q,C_q)$ равны
\[
\begin{array}{c|rrrrrrrrr}
q&-128&-127&-98&-97&-72&-71&-2&-1&1\\\hline
C_q&-200&12&3600&24&-200&12&400&-24&-24 .
\end{array}
\]
Они получены умножением конечных рядов исходного символа. Условие $\Im H(\kappa_m/N)=0$ даёт ненулевой рациональный многочлен степени 256 после умножения на $z^{128}$ суммы $\sum_q C_q(z^q-z^{-q})$. При $k\ge3$ первообразный корень $z=e_N(\kappa_m)$ имеет степень $6\cdot7^{k-1}>256$. Поэтому $F'_m(0)\ne0$ и $F_m(\tau N)$ трансцендентно. На общем нечётном модуле то же доказательство работает при $\varphi(N)>256$ и обратимости $\det(S^{2m}-I)$.

Если $3\mid m$ и $7\nmid m$, то $b$ из (4) равен единице. В $\Tr(A M^m)$ остаются лишь метки $(0,\pm7)$, обе с коэффициентом $9/20$. Запишем $b_m=7B$, $a_m-1=49A_0$, где $B,A_0$ единицы по модулю 7. Вторая делимость следует из $a_m^2-2b_m^2=1$ и обратимости $a_m+1$. Матрица $(S^m-I)/7$ имеет определитель $-2A_0$. При $v=7z$ формула (3) даёт фазу $e_{N/7}(-\bar4(B/A_0)Q(z))$. Подставляя $z=(0,\pm1)$, получаем $F'_m(0)\ne0$: синус угла $2\pi\bar2(B/A_0)/(N/7)$ ненулевой. Вместе два случая охватывают все $21\nmid m$. Для $o\mid m$ чистая степень скалярна, $F_m(0)=N^2$, тогда как $U_N(\tau N)^m$ нескалярна по теореме о строгом сжатии вейлевых вероятностей из части о жёсткости. Значит, $F_m(\tau N)<N^2$ и это число также трансцендентно. Здесь $N\ne7$, как требует указанная теорема.
\newpage
\sect{3. Оценка произведения на гейзенберговском времени}
Пусть $A_j=M^j A M^{-j}$, $c_A=\sum_d |c_d|\le6/5$. Сумму операторов внутри отрезка длины $l$ оценим в норме $\|X\|_{\rm n}=(N^{-1}\Tr X^*X)^{1/2}$. Для $1\le l\le m=3N/7$ каждая метка $S^j d$ в орбите данного $d$ встречается не более семи раз по (4). Поэтому
\[
 \left\|\sum_{j=j_0}^{j_0+l-1} A_j\right\|_{\rm n}
   \le c_A\sqrt{7l}<\tfrac{16}{5}\sqrt l .
\]
Сначала суммируем по одной метке с учётом кратностей в ортонормированном базисе Вейля, затем применяем треугольное неравенство по $d$. Граница равномерна по началу отрезка.

В одном блоке положим $V_j=\prod_{i=0}^{j-1}e^{itA_{j_0+i}}$ в порядке возрастания $i$, $S_j=\sum_{i=0}^{j-1}A_{j_0+i}$. Из телескопирования и $\|e^{it A}-I-it A\|_{\rm op}\le t^2\|A\|_{\rm op}^2/2$:
\[
 \|V_l-I\|_{\rm n}\le |t|\,\left\|\sum_{j=0}^{l-1}V_j A_{j_0+j}\right\|_{\rm n}
                             +\tfrac12 t^2 l c_A^2 .
\]
Дискретное суммирование по частям превращает сумму в $V_{l-1}S_l+\sum_{j=1}^{l-1}(V_{j-1}-V_j)S_j$. Норма каждого $V_{j-1}-V_j$ не больше $|t|c_A$, а $\|S_j\|_{\rm n}\le(16/5)\sqrt j$. Если $|t|l\le1$, получаем $\|V_l-I\|_{\rm n}\le8|t|\sqrt l$. При $0<|t|\le1/2$ разбиваем $m$ шагов на блоки длины $\min(m,\lfloor |t|^{-1}\rfloor)$ и возможный остаток. Если весь обход входит в один блок, применяем полученную оценку. Иначе число блоков не больше $2m/\lfloor |t|^{-1}\rfloor$, а $\lfloor |t|^{-1}\rfloor\ge 1/(2|t|)$. Отсюда
\[
 \|U_N(t)^m-M_N^m\|_{\rm n}
                       \le 8 |t|\sqrt m+23 m|t|^{3/2}.             \tag{7}
\]
Все множители унитарны, поэтому нормы блоков складываются телескопически. Для $t=0$ оценка очевидна.

Подставляя $t=Nu_N$, видим, что правая часть (7) стремится к нулю, если $u_N=o(N^{-5/3})$. При этом $M_N^m$ скалярна. Следовательно, $|\Tr U_N(t)^m|/N\to1$ на этих временах. Это исключает $\si$-сходимость в указанном вспомогательном режиме по необходимому условию слабого закона из работы о возврате Пелля. Прямой локальный тест на следующей странице даёт то же следствие.

Для исходного $\tau$ и $k\le56$ рациональная проверка квадратов двух членов (7) даёт их сумму $<0.204$. Достаточно проверить верхнюю размерность, поскольку правая часть возрастает при $t=\tau N$, $m=3N/7$. Пусть $M_N^m=\mu I$. Точное тождество
\[
 \Re\!\left(\bar\mu\,\frac{\Tr U_N^m}{N}\right)
             =1-\frac12\|U_N^m-M_N^m\|_{\rm n}^{\,2}                \tag{8}
\]
даёт нормированный след больше $0.979$. В частности, первая оценка (1) выполняется с запасом.
\newpage
\sect{4. Полный парный тест и локальный свидетель в конечном диапазоне}
Пишем $w(t)=(\phi*\phi)(t)^2/(4I^4)$ с $I=\int\phi$, как в исходной задаче. Функция $w$ не возрастает на $[0,2]$: симметричность и унимодальность свёртки следуют из разложения неотрицательных унимодальных функций на индикаторы симметричных интервалов. Рациональный расчёт в приложении даёт
\[
 w(3/7)>0.06854,\qquad w(0)<0.1141,\qquad
 C_w:=\int_0^2 w(t)\min(t,1)\,dt<0.02067.
\]
На интервале $[(m-1)/N,m/N]$ правая дискретная точка удовлетворяет $\min(m/N,1)w(m/N)\le(\min(t,1)+1/N)w(t)$. Поэтому эталонная часть (P) не больше $C_w+2w(0)/N$. Все члены квадрата следа неотрицательны. Оставив лишь $m=3N/7$ и применив (8), при $N\ge49$ получим
\[
 \Delta_{2,N}(b_*)>
         0.06854\,(1-0.204^2/2)^2-0.02067-\frac{0.2282}{49}>0.040.
\]
Это свидетель для полного заданного веса, с исходной силой кика и $m$ порядка размерности. Численный предел по $N$ здесь не используется.

\subt{Ограниченный непрерывный наблюдаемый} Пусть $T_0=70$ и $f(y)=(1-|y|/T_0)_+$. Для развёрнутой меры $\Xi_N$ с равномерным началом обозначим
\[
 X=\Xi_N(f)/T_0,\qquad Y=\Xi_N(e^{2\pi i (3/7)y}f)/T_0,\qquad
 G=\min(1,X-|Y|).
\]
Это непрерывный функционал со значениями в $[0,1]$. Имеем $\EE X=1$. Коэффициент Фурье случайной величины $Y$ по равномерному началу при частоте $m=3N/7$ имеет модуль $|\Tr U_N^m|/N$: после суммирования по периодам и интегрирования по началу остаётся $\int f/T_0=1$. Следовательно, $\EE G\le1-\EE|Y|\le1-|\Tr U_N^m|/N<0.204^2/2$ в заявленном диапазоне.

У процесса $\si$ дисперсия $\Xi(g)$ для комплексного $g\in L^2$ не больше $\|g\|_2^2$. Это следует из двухточечной функции и неотрицательности фурье-множителя $(1-|t|)_+$ квадрата ядра. Значит, $\EE_{\si}|X-1|\le\sqrt{2/(3T_0)}$. Среднее $Y$ равно нулю, поскольку преобразование Фурье $f$ пропорционально $\sinc^2(\pi T_0 t)$ и $T_0(3/7)=30$. Та же оценка дисперсии даёт $\EE_{\si}|Y|\le\sqrt{2/(3T_0)}$. Так как $G\ge1-|X-1|-|Y|$, получаем
\[
             \EE_{\si}G-\EE G
                  >1-2/\sqrt{105}-0.204^2/2>0.78 .
\]
Проверка относится к одному функционалу для всех $2\le k\le56$. Для убывающей силы $u_N=o(N^{-5/3})$ та же оценка с (7) даёт $\EE G\to0$, а значит, отвергает $\si$ на бесконечном ряду этого вспомогательного семейства.
\newpage
\sect{5. Точный предельный парный тест при слабом кике}
Дальнейшие асимптотики относятся к вспомогательному режиму
\[
          |t_N|\le\frac{\log N}{5N},\qquad N=7^k,\qquad
          L=\max_{m<2N}|m t_N|\le\tfrac25\log N.                    \tag{9}
\]
При $m<2N$ раскрываем $U_N(t)^m=\bigl(\prod_{j=0}^{m-1} e^{it A_j}\bigr) M^m$. Нормированный относительно $p_m(0)$ след каждой вейлевой метки имеет модуль не больше 1 по (3), откуда $|p_m(t)/p_m(0)|\le e^{c_A |mt|}$.

Пусть $3\mid m$ и $b=k-a$ из (4). Во всяком ненулевом члене тейлоровского разложения должно выполняться
\[
              \sum_{i=1}^r S^{j_i}d_i\equiv0\bmod 7^b,\qquad
                        0\le j_1\le\cdots\le j_r<m .
\]
Нулевая метка отсутствует в $A$. При фиксированных остальных временах первое время пробегает одну орбиту заданной метки $d_1$ по модулю $7^b$. Если $b>\nu_7(d_1)$, возможных $j_1$ не больше $m/(3\cdot7^{b-\nu_7(d_1)-1})+1$. При $b\le\nu_7(d_1)$ достаточно тривиальной границы $m$. В обоих случаях число не превосходит $C\,7^a m/N$, поскольку $m\ge3\cdot7^{b-1}$, $\nu_7(d_1)\le1$. Для повторных времён тейлоровские знаменатели нельзя просто убрать при подсчёте остальных времён. Если время $j_1$ имеет кратность $h$, то после удаления первого индекса остаётся коэффициент с множителем $1/(h-1)!$ вместо $1/h!$, то есть увеличивается. Сумма этих увеличенных весов по всем упорядоченным оставшимся временам точно $m^{r-1}/(r-1)!$. Действительно, каждое мультимножество соответствует $(r-1)!/\prod h_j!$ размеченным выборам времён. Следовательно,
\[
          \left|\frac{p_m(t_N)}{p_m(0)}-1\right|
             \le C\,\frac{7^a}{N} L e^{c_A L}
                      \qquad(3\mid m).                             \tag{10}
\]
Для $3\nmid m$ имеется $|p_m(0)|=1$. Из (10), в том числе при $b=k$, вытекает
\[
 \frac1{N^2}\sum_{1\le m<2N}|p_m(t_N)-p_m(0)|^2
                   \le C N^{-1}(1+L)^2 e^{2c_A L}=o(1).            \tag{11}
\]
Так как $2c_A L\le(24/25)\log N$, сходимость равномерна в (9).

По (4) чистая арифметическая сумма квадратов следов, нормированная на $N^2$, ограничена. Класс $b=k-a$ содержит $O(7^a)$ времён ниже $2N$ и даёт $O(7^{-a})$; времена $3\nmid m$ дают $O(N^{-1})$. Потому из (11) и Коши предел полного парного теста равен
\[
 \boxed{\displaystyle
      \sum_{a=0}^{\infty}7^{-2a}
            \sum_{\substack{u\ge1,\ 7\nmid u\\3u<2\cdot7^{a+1}}}
                  w(3u/7^{a+1})-C_w .}                            \tag{12}
\]
При $a=0$ все $u$ автоматически взаимно просты с 7. Ряд сходится геометрически. Уже первое слагаемое $w(3/7)$ превышает $C_w$. Для полного выражения рационально сертифицирован интервал $(0.084,0.087)$.
\newpage
\sect{6. Из временных следов к кратностям на решётке}
Пусть $\gamma_N$ --- общая фаза линейной решётки, то есть $M_N^o=e^{2\pi i o\gamma_N}I$, $o=3N/7$. Округлим фазу $\theta_i$ оператора $U_N(t_N)$ до ближайшего узла $\gamma_N+j_i/o$. Ошибку в единицах шага обозначим $\epsilon_i\in[-1/2,1/2]$. Взяв $m=o$ в (10), получаем
\[
 \frac1N\sum_i \epsilon_i^2
 \le\frac18\left(1-\Re\frac{p_o(t_N)}{p_o(0)}\right)=o(1),
\]
поскольку $1-\cos(2\pi x)\ge8x^2$ на указанном отрезке. Для равномерного начала развёртки это обеспечивает близость локальных тестов исходного и округлённого процессов. Для $g\in C_c^1$ ожидаемый абсолютный сдвиг теста не больше $\|g'\|_1\sum_i |\epsilon_i|/o=o(1)$, интегрированием по началу и Коши.

Сравним округлённые кратности с чистыми. Введём неотрицательное кардинальное ядро
\[
      K_o(x)=o^{-2}\left|\sum_{l=0}^{o-1} e^{2\pi i l x}\right|^2 .
\]
Сумма $K_o(x-j/o)$ по всем узлам равна 1, а в узлах ядро равно единице только в нуле. При $|\epsilon|\le1/2$ масса, уходящая из ближайшего узла, не больше $C\epsilon^2$. Это следует из $K_o(\epsilon/o)=(\sin(\pi\epsilon)/(o\sin(\pi\epsilon/o)))^2$ и $\sin y/y\ge1-y^2/6$ на $[-\pi/2,\pi/2]$. Следовательно, в норме $\ell^1$ сглаженные массы
\[
           s_j=\frac1N\sum_i K_o(\theta_i-\gamma_N-j/o)
\]
отличаются от округлённых масс $n_j/N$ на $o(1)$. Для спектра $M$ сглаживание ничего не меняет, так как все фазы лежат на узлах.

Фурье-коэффициенты ядра равны $o^{-1}(1-|m|/o)$ при $|m|<o$. В дискретном преобразовании по $j$ каждый индекс получает не более двух алиасов $m$ из $( -o,o )$. Нулевая разность следов тождественно равна нулю. Парсеваль и Коши дают
\[
      \sum_{j=0}^{o-1}|s_j-n_j(M)/N|
           \le \left(\frac2{N^2}\sum_{0<|m|<o}
                       |p_m(t_N)-p_m(0)|^2\right)^{1/2}=o(1)
\]
по (11). Тем самым $\sum_j |n_j-n_j(M)|/N=o(1)$. При общем равномерном начале развёртки разность интегралов локального теста между двумя сеточными процессами в среднем не больше этой величины, умноженной на $\|g\|_1$. Для доказательства сходимости точечных процессов достаточно взять счётное определяющее семейство непрерывных компактно поддержанных тестов. Равномерный первый момент массы компактов даёт тесноту. Так получаем тот же слабый предел, что у $M_N$.
\newpage
\sect{7. Полный локальный закон линейной решётки}
Над $\Z/7^k\Z$ существуют два корня $\pm\rho$ уравнения $\rho^2=2$ с $\rho\equiv3\bmod7$. Поднятие строится рекуррентно единственно, поскольку $2\rho$ обратимо. Тогда $S$ диагонализуема в симплектическом базисе со значениями $\gamma=3+2\rho$ и $\gamma^{-1}$. Нормируем один из двух векторов так, чтобы определитель матрицы перехода был равен единице.

У любой такой матрицы перехода $C$ есть унитарное представление в базисе Вейля. Действительно, отображение $T(v)\mapsto T(Cv)$ сохраняет точное произведение и сопряжение и потому является $*$-автоморфизмом полной матричной алгебры. Оно внутреннее: достаточно применить его к матричным единицам и выбрать согласованный ортонормированный базис их образов. Оператор диагонального действия $\operatorname{diag}(\gamma,\gamma^{-1})$ представлен перестановкой $\psi(x)\mapsto\psi(\gamma^{-1}x)$. Отсюда $M_N$ унитарно сопряжён с этой перестановкой с точностью до общей фазы.

Слой элементов аддитивного порядка $7^l$ имеет $6\cdot7^{l-1}$ элементов и разбивается ровно на два цикла длины $3\cdot7^{l-1}$ по (4). Плюс имеется неподвижный ноль. Поэтому на узле $j/o$, $0\le j<o$, кратность равна
\[
         n_j(M)=2\sum_{a=0}^{k-1}\mathbf1_{7^a\mid j}+\mathbf1_{j=0}.
                                                                     \tag{13}
\]
При равномерном начале на окружности индекс соседнего узла равномерен по $\Z/o\Z$, а положение внутри щели равномерно и независимо. В развёрнутых единицах щель равна $N/o=7/3$. Любая конечная группа индексов соседних узлов имеет в пределе кратности $2(1+\nu_7(B+j))$. Проверим это, сначала оставив только условия делимости до степени $7^D$. Тогда начальный индекс равномерен по $\Z/7^D\Z$ для $k>D$. Вероятность, что хотя бы в одном из конечного числа узлов нужен следующий разряд, не больше числа узлов, умноженного на $7^{-D}$. После $D\to\infty$ получается указанная совместная сходимость. Единственный дополнительный нуль у перестановки попадает в фиксированное локальное окно с вероятностью $O(N^{-1})$. Границы окна не создают трудности из-за непрерывного сдвига. Это доказывает локальный предел для $M_N$ и для всего режима (9) по разделу 6.

В частности, предельный процесс имеет кратные точки с вероятностью 1, а средняя кратность узла равна $2\sum_{a\ge0}7^{-a}=7/3$, поэтому интенсивность процесса равна 1. Синус-процесс почти наверное прост. Результат использует силу $u_N$ из (9), что соответствует $u_N\to0$. Для фиксированного $\tau$ утверждение (1) даёт конечный, но не бесконечный ряд.

\textbf{Роль обычного следа.} В доказательствах (7)--(13) нулевая вейлева метка учитывается непосредственно через $\Tr U_N(t)^m$. Именно след на времени $o$ прикрепляет фазы к решётке. Моменты остальных времён восстанавливают кратности после сглаживания ядром Фейера.
\newpage
\sect{8. Совместное распределение времени и фазы Кэли}
Вернёмся к временам $3\nmid m$. В обозначениях предыдущего раздела $\gamma=3+2\rho$. При $m=3s+r$, $r=1,2$, имеем $z=\gamma^r(\gamma^3)^s$ и
\[
                    \kappa_m=f(z)=-\frac{z+1}{4\rho(z-1)}
                             \pmod{7^k}.
\]
Положим $P=7^{k-1}$, $b=\lceil k/2\rceil-1$, $h=k-b-1$. Тогда $\gamma^3$ имеет порядок $P$, а $(\gamma^3)^{7^b}=1+7^{b+1}\beta$ с единицей $\beta\bmod7$. Разбиваем $s=s_0+7^b v$, $0\le s_0<7^b$, $0\le v<7^h$. Поскольку $2(b+1)\ge k$, разложение дроби по $v$ линейно в кольце $\Z/N\Z$. Для совместной моды
\[
      \frac1P\sum_{s\bmod P}
                  e_P(\ell s) e_N(n\kappa_{3s+r})
\]
внутренняя сумма по $v$ обращается в нуль, кроме условий
\[
        \ell+n\beta z f'(z)\equiv0\bmod 7^h,\qquad z=\gamma^r(\gamma^3)^{s_0}.
\]
На фиксированном классе $z\bmod7=2$ или 4 функция $z f'(z)=z/(2\rho(z-1)^2)$ имеет производную, обратимую по модулю 7: она пропорциональна $-(z+1)/(z-1)^3$. Поэтому в каждом таком классе условие задаёт не более одного класса по модулю $7^{h-\nu_7(n)}$, если $n\ne0$ и $h>\nu_7(n)$. Утверждение следует поднятием по одному разряду с обратимой производной. Поскольку $b\ge h-1$, эти классы посещаются последовательностью $s_0$ с равной частотой. Таким образом, совместная экспоненциальная сумма по модулю не больше $7^{1+\nu_7(n)-h}$ для $n\ne0$, равномерно по $\ell$. При $n=0$ и фиксированном $\ell\ne0$ она равна нулю.

Из этого следует совместная равномерность $(s/P,\kappa_{3s+r}/N)$ и, повторяя периоды, независимость предельной фазы Кэли от нормированного времени на любой фиксированной полосе $aN\le m\le b_0 N$. При отсечении по интервалу для $s$ второе утверждение следует либо из сходимости на прямоугольниках, либо из Фурье-завершения: $\ell^1$-норма коэффициентов индикатора интервала на $\Z/P$ имеет порядок $O(\log P)$ по сумме геометрической прогрессии. По (6) закон $F'_m(0)/m$ на временах $3\nmid m$ есть образ равномерной меры при функции $-2\Im H$. Его среднее нуль, дисперсия $412527/250000$. Число получено из суммы квадратов коэффициентов синуса. Таким образом, первый отклик обычного следа имеет ненулевую дисперсию на каждой гейзенберговской полосе.
\newpage
\sect{9. Многочастотные суммы для высших порядков}
Положим $aN\le m\le b_0N$, $3\nmid m$, и для каждой ненулевой метки $d_i$ её координаты в расщеплении обозначим $u_i=d_{i,1}+\rho d_{i,2}$, $v_i=d_{i,1}-\rho d_{i,2}$. Произведение $u_iv_i=Q(d_i)\ne0$ имеет ограниченный 7-адический порядок на конечном носителе. Для строго упорядоченных времён $j_1<\cdots<j_r$ положим $z_i=\gamma^{j_i}$. По произведению Вейля и формуле (5) фаза следа, делённого на $p_m(0)$, имеет вид
\[
 \kappa_m\sum_i Q(d_i)+
 \sum_{i<j}\left[
       A_m u_i v_j\,\frac{z_i}{z_j}
        +B_m v_i u_j\,\frac{z_j}{z_i}\right],                       \tag{14}
\]
где $A_m=\kappa_m-\bar4/\rho$, $B_m=\kappa_m+\bar4/\rho$. Обе константы обратимы: если $z=\gamma^m$, то $A_m=-z/(2\rho(z-1))$, $B_m=-1/(2\rho(z-1))$.

При $r\ge2$ рассмотрим полную нормированную экспоненциальную сумму по $s_1,\ldots,s_r\bmod P$ при фиксированных остатках $j_i\bmod3$, с произвольными добавочными модами $e_P(\sum_i \ell_i s_i)$. Разбиваем каждый $s_i$ на младшую часть и хвост по схеме раздела 8. Линейность хвостов снова точная, поскольку произведения двух хвостовых приращений равны нулю по модулю $7^k$. Сумма по первому хвосту нулевая, кроме условия на первую логарифмическую производную фазы. Если фиксированы все младшие части, кроме второй, оно имеет вид
\[
                 A z_2+B/z_2+C\equiv0\bmod 7^h.
\]
В частности, $A,B$ суть произведения единиц и компонент двух выбранных меток, и их 7-адические порядки ограничены абсолютной константой. Умножаем на $z_2$ и дополняем квадрат. Число квадратных корней произвольного вычета по модулю $7^h$ не больше $2\cdot7^{h/2}$: при ненулевом вычете чётного порядка $2v$ после деления остаётся два простых корня и $7^v$ поднятий, а для нуля число равно $7^{\lfloor h/2\rfloor}$. Аффинное отображение $z_2\mapsto2A z_2+C$ имеет ограниченное число прообразов. Так как $b\ge h-1$, в вычетах нужного класса младшая часть равномерна по модулю $7^h$. Доля решений есть $O(7^{-h/2})=O(N^{-1/4})$. Эта граница не зависит от $r$ и добавочных мод, поскольку оценка использует только вторую переменную и ограниченный носитель. Следовательно, все полные нормированные суммы имеют модуль $O(N^{-1/4})$.

Переносим это на упорядоченные времена. Разбиваем $j_i$ на номер полного периода $o$ (возможных номеров ограниченное число), остаток по модулю 3 и младший индекс $s_i$ в $[0,P)$. При равном номере периода условия $j_i<j_{i+1}$ дают сравнения $s_i<s_{i+1}$ или $s_i\le s_{i+1}$. Индикатор $x<y$ на $0\le x,y<P$ равен $\{(x-y)/P\}-x/P+y/P$. Каждый из трёх членов имеет фурье-норму $\ell^1=O(\log P)$ на циклической группе. Это видно из дискретной производной пилообразной функции и знаменателя $|1-e_P(j)|$. Для $\le$ добавляется диагональный индикатор с нормой 1. Отсечения каждого $j_i<m$ стоят ещё по $O(\log P)$ за индикатор интервала. Поэтому сумма (14) по строго упорядоченным временам, делённая на $m^r$, есть
\[
                        O\!\left(N^{-1/4}(C\log N)^{2r}\right)     \tag{15}
\]
с $C$ геометрической полосы, равномерно по $r$ и по меткам. Все суммы по частотам оценены по норме произведения свёрток.
\newpage
\sect{10. Комплексный след при растущем вспомогательном отклике}
Равномерно при $aN\le m\le b_0 N$, $3\nmid m$ и
\[
                     |mt|\le c_0\,\frac{\log N}{\log\log N}
\]
для достаточно малого фиксированного $c_0>0$ получена асимптотика
\[
 \boxed{\displaystyle
       \frac{p_m(t)}{p_m(0)}
             =1+i mt\,H(\kappa_m/N)
                    +O\!\left(e^{-c_1\log N/\log\log N}\right).}     \tag{16}
\]
Здесь $H$ определена в (6). Для доказательства раскроем произведение по порядкам до $R=\lfloor\log N/(100\log\log N)\rfloor$. В порядках $2\le r\le R$ вклад строго различных времён оценивается суммой
\[
           C N^{-1/4}\sum_{r=2}^R
                    [(c_A |mt|)(C\log N)^2]^r=o(N^{-1/8})
\]
по (15) и конечной $\ell^1$-норме символа. Если ряд по времени имеет повторы, его точные множители $1/\prod k_j!$ дают при фиксированном $r$ нормировку $m^r/r!$. Доля повторов при $r$ независимых выборах времён не больше $r(r-1)/(2m)$. Складывая все порядки, их общий абсолютный вклад не выше $(c_A |mt|)^2 e^{c_A|mt|}/(2m)=N^{-1+o(1)}$. Наконец, хвост по $r>R$ не превосходит хвоста экспоненты $\sum_{r>R}(c_A|mt|)^r/r!$, который не более $C2^{-R}$ при достаточно малом $c_0$. Члены нулевого и первого порядков точны. Так доказывается (16) без замены квантовых фаз вероятностями.

Например, при $t=\xi/N$ и фиксированном вещественном $\xi$ предельный эмпирический закон $|p_m(t)|^2$ на временах $3\nmid m$ есть закон $|1+i\xi V H(X)|^2$, где $V$ равномерен на рассматриваемой полосе нормированных времён и $X$ независимо равномерен на $[0,1]$. Это следует из (16), $|p_m(0)|=1$ и раздела 8. Разложение описывает настоящие модули следов на гейзенберговских временах вспомогательного семейства. На исходной силе $t=\tau N$ величина $|mt|$ имеет порядок $\tau N^2$ и выходит из этой области.

\textbf{Точные проверки.} В архиве независимо проверены коэффициенты $H$, фазовая формула Кэли, 7-адические классы, оценка (7) на верхнем конце заявленного диапазона, рациональные границы веса и полная арифметическая сумма (12). Матричные проверки на малых $7^k$ сверяют отношение вставленного вейлева следа с (3), а также фазу (14) для нескольких моментов. Неравенства (1) выведены аналитически, их константы проверены целыми числами.

\textbf{Линдеман--Вейерштрасс.} Применяется стандартная формулировка о линейной независимости $e^{\alpha_j}$ над $\overline{\mathbb Q}$ для различных алгебраических $\alpha_j$. См. \href{https://encyclopediaofmath.org/wiki/Lindemann_theorem}{Encyclopedia of Mathematics} и A. Baker, \emph{Transcendental Number Theory}, Cambridge University Press, гл.~1. Тождества Вейля и необходимые следствия слабого закона выведены в соответствующих предыдущих частях.
\newpage
\sect{11. Направленные границы для исходного веса}
Для воспроизводимости все границы, входящие в (1) и (12), выражены рациональными числами. Используем $D=7^3\cdot16=5488$ ячеек на единицу и масштаб $Q_0=2^{44}$ для целочисленных значений. В узлах $x=i/D$ внутри носителя функции $\phi$ вычисляем $v=(1-x^2)^{-1}$. Если $v\ge40$, берём границы $0$ и $Q_0^{-1}$: из $e>8/3$ следует $e^{-40}<Q_0^{-1}$. Если $v<40$, чередующиеся частичные суммы до степеней 15 и 16 обрамляют $e^{-v/64}$. Округляем их наружу на сетке $Q_0^{-1}\Z$ и шесть раз возводим в квадрат, каждый раз с направленным округлением. Так получаются точные интервалы для $e^{-v}$.

По монотонности $\phi$ на каждой половине носителя минимальные нижние и максимальные верхние оценки смежных узлов дают две ступенчатые границы. Пусть $l_i,h_i$ --- их целые числители. Тогда $I$ лежит между $\sum l_i/(D Q_0)$ и $\sum h_i/(D Q_0)$. Свёртки этих двух ступенчатых функций в $j/D$ равны соответственно $\ell_j/(D Q_0^2)$ и $u_j/(D Q_0^2)$, где $\ell_j,u_j$ вычисляются обычным целочисленным умножением полиномов их числителей. Точное целое упаковочное умножение допустимо, если каждый разряд шире максимального коэффициента свёртки. Отсюда
\[
   \frac{D^2\ell_j^2}{4(\sum h_i)^4}
                  \le w(j/D)\le
   \frac{D^2u_j^2}{4(\sum l_i)^4}.
\]
Для $C_w$ умножаем на $\min(t,1)$ и берём правую границу для растущего множителя, левую для убывающего $w$ в верхней сумме Дарбу. В нижней сумме меняем границы. Так получены $w(3/7)>0.06854$, $C_w<0.02067$, $w(0)<0.1141$. Гарантированная нижняя граница (1), с использованием самих рациональных интервалов, больше $0.04040$.

В (12) первые три класса $a=0,1,2$ попадают на узлы выбранной сетки. У оставшегося хвоста верхняя оценка
\[
            w(0)\sum_{a\ge3} 7^{-2a}\,\frac{2\cdot7^{a+1}}3
                         =\frac{49}{9\cdot7^3}w(0).
\]
Все члены хвоста положительны, поэтому его можно опустить в нижней границе. С учётом направленных оценок для $C_w$ получен заключающий интервал $(0.08443,0.08656)$.

При $N=7^{56}$ и $t=N/10^{80}$ квадраты слагаемых в (7) меньше $(10^{-8})^2$ и $(203/1000)^2$ соответственно. Проверяются рациональные величины $64t^2m$ и $23^2m^2t^3$, $m=3N/7$. На всех меньших размерностях выражения убывают. Этим полностью сертифицированы константы возврата в (1). Численная сверка фаз на малых матрицах лишь проверяет конвенцию знаков и не входит в границы.
\end{document}
